\subsection{Standalone Hero Models}\label{sec:hero_models}
We report standalone ``hero'' experiments for STANet \cite{chen2020stanet} and BIT \cite{chen2022bit} under their native/custom pipelines. The goal is to test transferability of \cosa{} when inserted into strong non-baseline backbones.
Table~\ref{tab:hero_models} reports the changed-class results for these two backbones and shows that \cosa{} delivers a large gain on STANet but only a marginal change on BIT.

\begin{table*}[t]
  \centering
  \caption{Standalone hero models under their native pipelines (changed class).}
  \label{tab:hero_models}
  \setlength{\tabcolsep}{7pt}
  \renewcommand{\arraystretch}{1.05}
  \small
  \begin{tabular*}{\textwidth}{@{\extracolsep{\fill}} l l r r r r}
    \toprule
    Model & Variant & \PrecC{} & \RecC{} & \Fonec{} & \Iouc{} \\
    \midrule
    STANet & Baseline & 70.31 & \textbf{94.71} & 80.71 & 67.66 \\
    STANet & + \cosa{} & \textbf{87.11} & 85.19 & \textbf{86.14} & \textbf{75.66} \\
    BIT & Baseline & 91.37 & \textbf{91.08} & 91.22 & 83.87 \\
    BIT & + \cosa{} & \textbf{92.38} & 90.17 & \textbf{91.26} & \textbf{83.93} \\
    \bottomrule
  \end{tabular*}
\end{table*}


The two hero models exhibit clearly different responses. On STANet, \cosa{} provides a large gain (\Fonec{}: 80.71 \textrightarrow 86.14; \Iouc{}: 67.66 \textrightarrow 75.66) with a strong precision increase, indicating effective suppression of over-detection. On BIT, the effect is near-neutral (\Fonec{}: 91.22 \textrightarrow 91.26), consistent with partial redundancy between explicit correlation gating and BIT's transformer-based temporal interaction. Standalone evidence therefore matches the broader pattern in this paper: \cosa{} is most beneficial when it complements, rather than duplicates, the temporal interaction already encoded by the backbone \cite{w2022,wei2024}.
